\section{Lecture 8 (January 31, 2025, est.): Tensor functors and moderate growth}\label{lec8}
\noindent\textit{Scribe: Guillermo Sanmarco}


\subsection*{The parity element}
We first make the condition on the element $z$ from Lecture~\ref{lec7} explicit. Put $A=\mathcal O(G)$ and view $z\in G(k)$ as an even algebra homomorphism $z:A\to k$. The pullback of conjugation by $z$ is
\[
 a\longmapsto\sum z(a_{(1)})\,a_{(2)}\,z^{-1}(a_{(3)}).
\]
Here $z^{-1}=z\circ S$, and $z^2=1$ makes $z^{-1}=z$. The parity condition requires this expression to equal $a_0-a_1$ for $a=a_0+a_1$, $a_i\in A_i$. It does not require $z$ to be central in the whole supergroup.

\subsection*{Monoidal and symmetric tensor functors}
\begin{definition}
A \emph{strong monoidal functor} $(F,F_0,F_2):\mathcal C\to\mathcal D$ consists of a functor, an isomorphism $F_0:\mathbb{1}_{\mathcal D}\to F(\mathbb{1}_{\mathcal C})$, and natural isomorphisms
\[
 F_{2,X,Y}:F(X)\otimes F(Y)\longrightarrow F(X\otimes Y).
\]
They satisfy the associativity identity
\[
\begin{split}
 F(a_{X,Y,Z})F_{2,X\otimes Y,Z}(F_{2,X,Y}\otimes\mathrm{id})
 ={}&F_{2,X,Y\otimes Z}(\mathrm{id}\otimes F_{2,Y,Z})\\
 &\hspace{-10mm}\circ a_{F(X),F(Y),F(Z)},
\end{split}
\]
and the two unit identities
\[
\begin{split}
 F(l_X)F_{2,\mathbb{1},X}(F_0\otimes\mathrm{id})&=l_{F(X)},\\
 F(r_X)F_{2,X,\mathbb{1}}(\mathrm{id}\otimes F_0)&=r_{F(X)}.
\end{split}
\]
For braided categories, $F$ is \emph{braided} if
\[
 F(c^{\mathcal C}_{X,Y})F_{2,X,Y}
 =F_{2,Y,X}c^{\mathcal D}_{F(X),F(Y)}.
\]
If both braidings are symmetric, we call it a symmetric monoidal functor.
\end{definition}
In these notes, a \emph{symmetric tensor functor} between symmetric tensor categories means a $k$-linear exact strong symmetric monoidal functor. Such a functor is faithful.

\begin{proof}[Proof of faithfulness]
If $X\ne0$, the evaluation $X^*\otimes X\to\mathbb{1}$ is nonzero by the triangle identity, and hence surjective because the unit is simple. Exactness and $F(\mathbb{1})\simeq\mathbb{1}$ imply $F(X)\ne0$. If $f$ is a nonzero morphism, its image is nonzero. Exactness gives $\operatorname{im}F(f)\simeq F(\operatorname{im}f)\ne0$, so $F(f)\ne0$.
\end{proof}

\begin{lemma}\label{lem:trace-functor}
A symmetric tensor functor preserves categorical traces and dimensions.
\end{lemma}
\begin{proof}
Applying $F$ to evaluation and coevaluation, and inserting $F_0,F_2$, makes $F(X^*)$ a dual of $F(X)$. The triangle identities follow by applying $F$ to those for $X$. Symmetry compatibility identifies the image of the braiding in the trace formula with the braiding of $F(X)$ and $F(X^*)$. Thus
\[
 \operatorname{Tr}(F(f))=F_0^{-1}F(\operatorname{Tr}(f))F_0.
\]
On the scalar endomorphisms of the unit this is the same scalar because $F$ is $k$-linear. Set $f=\mathrm{id}_X$ for the dimension statement.
\end{proof}

For $\operatorname{char}k\ne2$, there is no symmetric tensor functor from $\mathrm{sVec}$ to $\operatorname{Rep}_k(G)$ for an ordinary affine group scheme $G$. Indeed, the odd line $\Pi$ has tensor square $\mathbb{1}$ and categorical dimension $-1$. Its image would have one-dimensional underlying space, hence dimension $1$, contradicting Lemma~\ref{lem:trace-functor}. The tensor-square condition is essential here: in positive characteristic an arbitrary vector space can have dimension congruent to $-1$ in $k$.

\subsection*{Moderate growth}
Write $\ell(X)$ for the composition length of an object in a finite-length abelian category.
\begin{definition}
An object $X$ has \emph{moderate growth} if there is an integer $d_X\geq1$ such that
\[
 \ell(X^{\otimes n})\leq d_X^n\qquad(n\geq0).
\]
A tensor category has moderate growth if every object does.
\end{definition}

\begin{lemma}
For an exact faithful functor $F:\mathcal C\to\mathcal D$ between finite-length abelian categories, $\ell(X)\leq\ell(F(X))$.
\end{lemma}
\begin{proof}
Apply $F$ to a composition series of $X$. Exactness preserves its successive quotients, and faithfulness makes their images nonzero. Each image has length at least one; additivity of length gives the inequality.
\end{proof}
Consequently, moderate growth pulls back along tensor functors. In $\operatorname{Rep}(G,z)$ the ordinary total vector-space dimension, rather than superdimension, gives
\[
 \ell(X^{\otimes n})\leq(\dim_kX_0+\dim_kX_1)^n.
\]

Categories with finitely many simple isomorphism classes also have moderate growth. To see this, choose representatives $S_1,\ldots,S_r$. For fixed $X$ put
\[
 d_X=\max\{1,\ell(X\otimes S_1),\ldots,\ell(X\otimes S_r)\}.
\]
Tensor exactness and a composition series of $Y$ give $\ell(X\otimes Y)\leq d_X\ell(Y)$. Starting with $Y=\mathbb{1}$ proves the required bound by induction. In a finite tensor category one can alternatively use the Frobenius--Perron dimension: its positivity on simples gives $\ell(X^{\otimes n})\leq\operatorname{FPdim}(X)^n$.

\begin{theorem}[Deligne]\label{thm:deligne-moderate}
Let $k$ be algebraically closed of characteristic zero and let $\mathcal C$ be a symmetric tensor category. The following are equivalent:
\begin{enumerate}
\item $\mathcal C$ has moderate growth;
\item there is a symmetric tensor functor $\mathcal C\to\mathrm{sVec}_k$;
\item $\mathcal C\simeq\operatorname{Rep}(G,z)$ as a symmetric tensor category for an affine supergroup $G$ and an element $z$ with $z^2=1$ and $\operatorname{Ad}(z)$ equal to parity.
\end{enumerate}
\end{theorem}
The implication from moderate growth to a super fiber functor is a substantial theorem, not proved here; see \cite[Theorems 9.11.4 and 9.11.9]{EGNO}. The length estimate above proves the converse. The reconstruction behind the remaining equivalence motivates our next lectures, starting with ordinary fiber functors to vector spaces.
